\section{Thuật toán A*}
\subsection{Mô hình hoá bài toán}
Nhóm chúng em áp dụng kiến thức môn Cơ Sở Trí Tuệ Nhân Tạo vào để tiến hành mô hình hoá bài toán trước khi cài đặt thuật toán, cụ thể như sau:
\begin{itemize}
    \item \textbf{State:} Mỗi trạng thái được biểu diễn bởi đối tượng
    \texttt{GameState}, gồm trạng thái khối và tập các cầu đang mở. Có thể
    ký hiệu trạng thái dưới dạng
    \[
        s = (P, o, a, B),
    \]
    trong đó $P$ là bộ tọa độ các ô mà khối chiếm giữ; $o$ là hướng của
    khối, nhận một trong bốn giá trị \texttt{STANDING} (đứng),
    \texttt{HORIZONTAL} (nằm theo trục $x$), \texttt{VERTICAL} (nằm theo
    trục $y$) hoặc \texttt{SPLIT} (tách thành hai khối con); $a$ là chỉ số
    khối con đang được điều khiển khi khối bị tách; $B$ là tập mã định danh
    của các cầu đang mở, được lưu bằng \texttt{frozenset}.

    Khi khối đứng, $P$ có một tọa độ; khi khối nằm hoặc bị tách, $P$ có
    hai tọa độ. Các tọa độ được sắp xếp để chuẩn hóa biểu diễn; khi khối
    bị tách, chỉ số $a$ được điều chỉnh theo thứ tự này để giữ đúng khối
    con đang được chọn. Trạng thái ban đầu $s_0$ được lấy từ
    \texttt{board.game\_state}. Hai trạng thái có cùng vị trí và hướng
    nhưng khác tập cầu đang mở hoặc khác khối con đang được chọn vẫn là
    hai trạng thái khác nhau.

    \item \textbf{Goal:} Trạng thái đích là một trạng thái hợp lệ mà khối
    nguyên vẹn đang đứng trên ô có kiểu \texttt{Goal}. Điều kiện này được
    kiểm tra bởi \texttt{GameRules.is\_goal}. Khối nằm ngang qua ô đích
    hoặc một khối con nằm trên ô đích chưa thỏa mãn điều kiện thắng.
    Mục tiêu của thuật toán là tìm chuỗi hành động hợp lệ từ $s_0$ đến
    một trạng thái đích với tổng chi phí nhỏ nhất.

    \item \textbf{Actions:} Tập hành động gồm \texttt{UP}, \texttt{DOWN},
    \texttt{LEFT}, \texttt{RIGHT} và \texttt{SWITCH}. Bốn hành động đầu
    làm khối nguyên vẹn lăn theo hướng tương ứng; nếu khối đã bị tách,
    chúng chỉ di chuyển khối con đang được chọn một ô. Hành động
    \texttt{SWITCH} chỉ hợp lệ khi khối bị tách và dùng để đổi khối con
    đang được điều khiển, không làm thay đổi vị trí.
    Tại mỗi trạng thái, \texttt{get\_valid\_actions} chỉ trả về các hành
    động hợp lệ. Mỗi hành động, kể cả \texttt{SWITCH}, có chi phí bằng
    $1$, nên chi phí đường đi $g(s)$ chính là số hành động đã thực hiện.

    \item \textbf{Transition Model:} Hàm \mbox{\texttt{apply\_action}}
    của \texttt{GameRules} xác định trạng thái kế tiếp:
    \[
        T(s, u) = s', \qquad u \in A(s),
    \]
    với $A(s)$ là tập hành động hợp lệ tại $s$. Đối với hành động di
    chuyển, hàm tính vị trí và hướng mới theo tư thế hiện tại, sau đó
    kiểm tra hình dạng khối và khả năng nâng đỡ của các ô bên dưới.
    Di chuyển ra ngoài bàn cờ, đi lên cầu đang đóng hoặc đứng trên ô
    dễ vỡ đều không hợp lệ.

    Sau khi khối đáp xuống, các công tắc được kích hoạt sẽ được xử lý
    theo thứ tự mã định danh để cập nhật cầu bằng hiệu ứng
    \texttt{TOGGLE} hoặc \texttt{SET}. Công tắc tách khối còn đưa hai
    khối con đến các vị trí được cấu hình. Trạng thái sau hiệu ứng được
    kiểm tra lại; nếu khối trước hành động đã ở dạng tách và hai khối
    con sau di chuyển nằm kề nhau theo cạnh, chúng được ghép thành một
    khối nằm. Riêng \texttt{SWITCH} chỉ đổi chỉ số khối con đang được
    chọn, không kích hoạt công tắc hay ghép khối.
    Hàm trả về trạng thái mới khi hợp lệ và \texttt{None} nếu hành động
    không hợp lệ, đồng thời không thay đổi trạng thái đầu vào.

    \item \textbf{Heuristic Function:} Thuật toán đánh giá mỗi nút bằng
    \[
        f(s) = g(s) + h(s).
    \]
    Trong cài đặt, $h(s)$ được tra từ \texttt{heuristic\_table}, được
    xây dựng một lần khi khởi tạo bộ giải bằng hàm
    \texttt{\_build\_heuristic\_table}. Nhóm sử dụng khoảng cách ngắn
    nhất trên một bài toán nới lỏng: tất cả cầu luôn mở và các hiệu ứng
    thay đổi cầu của công tắc được loại bỏ. Các quy tắc lăn khối, ô dễ
    vỡ, tách khối, dịch chuyển khối con, đổi khối con và ghép khối vẫn
    được giữ nguyên.

    Bộ giải liệt kê các cấu hình khối đứng, nằm và tách trên các vị trí
    của bàn cờ, bao gồm cả vị trí cầu. Với cấu hình tách, cả hai lựa
    chọn khối con đang được điều khiển đều được xét. Từ các cấu hình
    hợp lệ, bộ giải sinh các chuyển trạng thái và lưu danh sách tiền
    nhiệm của từng cấu hình. Sau đó, BFS ngược được thực hiện đồng
    thời từ tất cả cấu hình đứng trên ô đích để tính
    \[
        h(s) = d_{\mathrm{relaxed}}\bigl(\operatorname{block}(s),
        \mathcal{G}_{\mathrm{relaxed}}\bigr),
    \]
    trong đó $d_{\mathrm{relaxed}}$ là số hành động ít nhất trên đồ thị
    nới lỏng và $\mathcal{G}_{\mathrm{relaxed}}$ là tập cấu hình đích.
    Việc lưu tiền nhiệm giúp BFS đi ngược đúng các cạnh có hướng, kể
    cả khi công tắc gây dịch chuyển khối. Khóa tra cứu chỉ là
    \texttt{state.block}, không bao gồm tập cầu đang mở. Cấu hình đích
    có giá trị $0$; cấu hình không có đường đến đích trong bài toán
    nới lỏng nhận giá trị $+\infty$.

    Do việc mở sẵn mọi cầu và bỏ hiệu ứng thay đổi cầu không loại bỏ
    đường đi hợp lệ nào của bài toán gốc khi chỉ xét cấu hình khối,
    heuristic không đánh giá quá cao chi phí còn lại, tức là
    $h(s) \leq h^*(s)$. Với mỗi chuyển trạng thái hợp lệ $s \to s'$, ta
    còn có $h(s) \leq 1 + h(s')$. Vì vậy, heuristic có tính chấp nhận
    được và nhất quán, giúp A* tìm lời giải có số hành động ít nhất.

    \item \textbf{Avoid revisiting states:} Bộ giải sử dụng từ điển
    \texttt{g\_score} để lưu chi phí nhỏ nhất đã tìm được đến từng
    \texttt{GameState}. Ban đầu, $g\_score[s_0] = 0$. Khi sinh trạng
    thái kế tiếp $s'$ từ $s$, thuật toán tính $g_{\mathrm{new}} = g(s)+1$.
    Nếu $g_{\mathrm{new}} \geq g\_score[s']$, trạng thái này bị bỏ qua;
    nếu nhỏ hơn, thuật toán cập nhật chi phí và đưa nút mới vào hàng
    đợi ưu tiên. Trạng thái chưa có trong từ điển được xem là có chi
    phí $+\infty$.

    Hàng đợi \texttt{PriorityQueue} lưu bộ ba
    $(f(s),\,\texttt{counter},\,\texttt{node})$, ưu tiên nút có $f(s)$
    nhỏ nhất; bộ đếm tăng dần phân xử các nút có cùng giá trị $f$ theo
    thứ tự đưa vào. Khi lấy một nút ra, nếu chi phí của nút lớn hơn
    giá trị hiện có trong \texttt{g\_score}, nút đó bị bỏ qua vì đã có
    đường đi tốt hơn. Cơ chế này hạn chế duyệt lặp nhưng vẫn cho phép
    xét lại một trạng thái khi tìm được đường đi có chi phí thấp hơn.
    Khi lấy được nút đích, bộ giải truy vết các liên kết cha để trả về
    chuỗi hành động; nếu hàng đợi rỗng mà chưa gặp đích, bộ giải trả
    về \texttt{None}.
\end{itemize}


\subsection{Về hàm Heuristic}
Ban đầu, em cân nhắc thiết kế hàm heuristic dựa trên khoảng cách
Manhattan. Tuy nhiên, khoảng cách giữa các ô chưa phản ánh đầy đủ số
hành động trong Bloxorz: một lần lăn có thể làm tọa độ khối thay đổi
nhiều hơn một ô, tư thế khối ảnh hưởng đến cách di chuyển, và công tắc
tách khối có thể đưa các khối con đến vị trí khác trên bản đồ. Vì vậy,
nếu dùng trực tiếp khoảng cách Manhattan làm số bước còn lại thì chưa
thể bảo đảm tính chấp nhận được. Sau khi trao đổi, nhóm em lựa chọn
heuristic dựa trên khoảng cách ngắn nhất trong một bài toán nới lỏng,
được tính trước bằng \textbf{BFS ngược đa nguồn}.

\subsubsection{Xây dựng bài toán nới lỏng}
Bước đầu tiên là tạo một bản sao của bàn cờ, mở sẵn tất cả các cầu và
loại bỏ các hiệu ứng thay đổi cầu của công tắc. Trong mã nguồn, các
công tắc được sao chép với thuộc tính \texttt{effects} rỗng; tập
\texttt{open\_bridges} chứa mã định danh của toàn bộ cầu. Nhờ đó, việc
đi qua một cầu không còn phụ thuộc vào lịch sử kích hoạt công tắc.

Việc nới lỏng chỉ áp dụng cho cơ chế đóng, mở cầu. Các ô dễ vỡ vẫn
không cho phép khối đứng lên; công tắc tách khối vẫn tách và dịch
chuyển hai khối con đến vị trí đã cấu hình. Các quy tắc lăn, di chuyển
khối con, đổi khối con đang được điều khiển và ghép khối vẫn được
thực hiện bởi \texttt{GameRules}. Giữ lại các quy tắc này giúp giá trị
ước lượng phản ánh tư thế và khả năng di chuyển của khối.

Vì trạng thái cầu đã được cố định, mỗi đỉnh trong đồ thị nới lỏng chỉ
cần lưu một \texttt{BlockState}, gồm vị trí, hướng và chỉ số khối con
đang được chọn. Gọi đồ thị này là
$G_R = (V_R, E_R)$ và tập cấu hình đứng trên ô đích là
$\mathcal{G}_R$. Với trạng thái thực $s$, heuristic được định nghĩa bởi
\[
    h(s) = \min_{q_g \in \mathcal{G}_R}
    d_{G_R}\bigl(\operatorname{block}(s), q_g\bigr),
\]
trong đó $d_{G_R}$ là số hành động ít nhất giữa hai cấu hình trên đồ
thị nới lỏng. Đây là khoảng cách chính xác trong bài toán nới lỏng và
là ước lượng chi phí còn lại trong bài toán gốc.

\subsubsection{Xây dựng bảng khoảng cách bằng BFS ngược}
Hàm \texttt{\_build\_heuristic\_table} thực hiện các bước sau:
\begin{enumerate}
    \item Thu thập tập tọa độ của các ô và các cầu. Tại mỗi tọa độ,
    sinh cấu hình đứng; nếu có ô kề theo trục $x$ hoặc $y$, sinh thêm
    cấu hình nằm tương ứng. Với mỗi cặp tọa độ phân biệt, sinh cấu
    hình tách cho cả hai lựa chọn khối con đang được điều khiển.

    \item Gắn từng cấu hình với tập cầu luôn mở và kiểm tra bằng
    \texttt{is\_valid}. Với mỗi cấu hình hợp lệ $q$, thử cả năm hành
    động bằng \texttt{apply\_action}. Nếu một hành động đưa $q$ đến
    $q'$, thêm $q$ vào danh sách tiền nhiệm của $q'$. Như vậy, cạnh
    $q \to q'$ của đồ thị thuận được lưu để có thể duyệt ngược từ
    $q'$ về $q$.

    \item Gán khoảng cách $0$ cho tất cả cấu hình đứng trên ô đích và
    đưa chúng vào hàng đợi \texttt{deque}. Đây là các nguồn khởi đầu
    của BFS trên đồ thị đảo cạnh.

    \item Lần lượt lấy cấu hình $q'$ ở đầu hàng đợi. Với mỗi tiền
    nhiệm $q$ chưa có khoảng cách, gán
    \[
        \operatorname{distance}[q]
        = \operatorname{distance}[q'] + 1
    \]
    rồi thêm $q$ vào cuối hàng đợi. Mỗi cấu hình chỉ được gán khoảng
    cách ở lần đầu tiên được phát hiện. Khi hàng đợi rỗng, trả về từ
    điển khoảng cách làm bảng heuristic.
\end{enumerate}

Mọi hành động đều có chi phí $1$, kể cả hành động đổi khối con, nên
BFS duyệt theo từng lớp khoảng cách và lần phát hiện đầu tiên cho
khoảng cách ngắn nhất. DFS không bảo đảm tính chất này. Việc sử dụng
danh sách tiền nhiệm cũng cần thiết vì các chuyển trạng thái có thể
có hướng: không thể mặc định rằng thực hiện hành động ngược sẽ đảo
được hiệu ứng tách khối hoặc dịch chuyển do công tắc gây ra.

Trong quá trình chạy A*, hàm \texttt{\_heuristic\_function} chỉ tra
bảng bằng khóa \texttt{state.block}. Hai trạng thái thực có cùng cấu
hình khối nhưng khác tập cầu đang mở sẽ nhận cùng giá trị heuristic,
mặc dù chúng vẫn được A* phân biệt khi lưu \texttt{g\_score}. Cấu hình
không có trong bảng nhận giá trị \texttt{math.inf}, tức là không có
đường đến đích ngay cả trong bài toán nới lỏng. Giá trị này không làm
thuật toán kết thúc sớm: cài đặt hiện tại vẫn đưa nút vào hàng đợi
ưu tiên với giá trị $f$ tương ứng.

\subsubsection{Admissible and Consitent}
Xét một lời giải của bài toán gốc có $k$ hành động. Khi mở sẵn tất cả
cầu và bỏ các hiệu ứng thay đổi cầu, chuỗi cấu hình khối của lời giải
đó vẫn tạo thành một đường đi hợp lệ đến đích trong đồ thị nới lỏng.
Do BFS trả về khoảng cách ngắn nhất trên đồ thị này, ta có
\[
    h(s) \leq k.
\]
Chọn $k = h^*(s)$, với $h^*(s)$ là chi phí lời giải tối ưu từ $s$,
suy ra $h(s) \leq h^*(s)$. Vì vậy, heuristic có tính chấp nhận được.
Đặc biệt, trạng thái đích có $h(s)=0$; nếu $h(s)=+\infty$ thì trạng
thái thực cũng không thể có lời giải.

Đối với tính nhất quán, xét một chuyển trạng thái hợp lệ $s \to s'$.
Chuyển động tương ứng giữa hai cấu hình khối cũng tồn tại trong đồ
thị nới lỏng với chi phí $1$. Đi từ cấu hình của $s$ sang cấu hình
của $s'$, rồi theo đường ngắn nhất đến đích từ đó, là một phương án
đến đích. Khoảng cách ngắn nhất từ cấu hình của $s$ không thể lớn hơn
chi phí của phương án này, nên
\[
    h(s) \leq 1 + h(s').
\]
Do đó, heuristic có tính nhất quán. Kết hợp với chi phí hành động
bằng $1$ và cơ chế cập nhật chi phí tốt nhất trong \texttt{g\_score},
A* trả về lời giải có số hành động ít nhất khi tìm được đích.

\subsubsection{Chi phí tính toán và giới hạn}
Bảng heuristic được xây dựng một lần khi khởi tạo một đối tượng
\texttt{AStarSolver} và được dùng lại trong quá trình tìm kiếm của
đối tượng đó. Mỗi lần tra cứu có chi phí trung bình $O(1)$ nhờ sử
dụng từ điển, thay vì phải chạy một phép tìm đường mới cho từng nút.

Đổi lại, bước tiền xử lý phải liệt kê cả các cấu hình tách khối.
Nếu tập tọa độ có $n$ phần tử, riêng phần này có
$2\binom{n}{2}=n(n-1)$ cấu hình ứng viên, nên tổng số cấu hình ứng
viên tăng theo $O(n^2)$, kể cả khi bản đồ không có công tắc tách khối.
Sau khi có đồ thị, BFS ngược cần thời gian
$O(|V_R|+|E_R|)$; việc lưu đồ thị tiền nhiệm và bảng khoảng cách cần
bộ nhớ $O(|V_R|+|E_R|)$. Chi phí dựng đồ thị còn bao gồm việc kiểm tra
luật và tra cứu ô trên bàn cờ, nên không thể đồng nhất toàn bộ thời
gian tiền xử lý với riêng chi phí BFS.

Heuristic vẫn có thể nhỏ hơn đáng kể so với chi phí thực tế nếu lời
giải phải đi vòng để mở cầu. Đây là phần thông tin đã được bỏ qua
khi nới lỏng bài toán. Các kiểm thử hiện có đối chiếu độ dài lời giải
A* với BFS trên năm bản đồ \texttt{map01}, \texttt{map02},
\texttt{map04}, \texttt{map08} và \texttt{map09}; đồng thời kiểm tra
bất đẳng thức nhất quán trên các chuyển trạng thái được BFS khảo sát
và tính chấp nhận được dọc theo lời giải tối ưu tìm được.
